\chapter{TIA}
\index{TIA}

\section*{Definition and Overview}
A Transient Ischaemic Attack (TIA) is defined as a brief episode of neurological dysfunction caused by focal brain, spinal cord, or retinal ischaemia, without acute infarction. Colloquially termed a ``mini-stroke'' or ``warning stroke,'' a TIA represents a critical clinical window that demands immediate medical evaluation. Historically, the diagnosis of a TIA relied on a time-based threshold, where neurological symptoms had to resolve completely within 24 hours. However, the advent of modern diagnostic neuroimaging, particularly diffusion-weighted magnetic resonance imaging (DWI-MRI), has revolutionized this paradigm. Today, clinical guidelines endorse a tissue-based definition. This shift recognizes that even brief episodes of ischaemia can result in permanent brain tissue injury (infarction). If neuroimaging demonstrates acute ischaemic tissue damage, the event is classified as an ischaemic stroke, regardless of symptom duration. Conversely, if symptoms resolve and there is no radiographic evidence of acute infarction, the diagnosis of TIA is confirmed.

The clinical significance of a TIA lies in its role as a potent, immediate precursor to a future, potentially devastating ischaemic stroke. A patient who has experienced a TIA is at extremely high risk of experiencing a completed stroke in the subsequent hours, days, and weeks. This risk is front-loaded, with the highest probability of stroke occurring within the first 24 to 48 hours. Identifying, evaluating, and treating a TIA rapidly is therefore one of the most effective preventive interventions in cardiovascular medicine. Prompt and aggressive management can reduce the subsequent risk of a major stroke by up to 80\%, saving lives and preventing long-term physical, cognitive, and emotional disability. Understanding TIA from a clinical and public health perspective is essential for reducing the global burden of cerebrovascular disease.

\section*{Epidemiology}
The epidemiology of Transient Ischaemic Attacks provides vital insight into the demographic groups most at risk and the overall burden of cerebrovascular disease. Determining the precise incidence and prevalence of TIA is inherently challenging because many individuals experience transient symptoms---such as a brief spell of slurred speech or temporary numbness---and fail to seek medical attention, dismissing the episode as a minor occurrence. Consequently, epidemiological data likely underestimate the true burden of the disease. In developed nations, the annual incidence of TIA is estimated to range between 50 and 100 cases per 100,000 people. This translates to hundreds of thousands of new cases each year globally.

The incidence of TIA is heavily influenced by age. It is relatively rare in individuals under the age of 50, but the incidence rises exponentially thereafter. The majority of TIA cases occur in individuals aged 65 and older, reflecting the cumulative impact of chronic cardiovascular risk factors and age-related vascular changes. In terms of sex, men have a higher age-adjusted incidence of TIA compared to women, particularly in middle-aged and early-elderly cohorts. However, as life expectancy increases, the absolute number of cases in older women rises, narrowing the sex gap.

Significant racial and ethnic disparities exist in TIA epidemiology. For example, African American and Hispanic populations in the United States, as well as Indigenous populations in many countries, Canada, and New Zealand, experience higher rates of TIA compared to their Caucasian counterparts. These disparities are closely linked to a higher prevalence of predisposing risk factors, including chronic hypertension, type 2 diabetes mellitus, obesity, and limited access to primary preventive healthcare.

A history of TIA is one of the strongest independent risk factors for stroke. Epidemiological studies show that approximately 15\% to 30\% of all patients who suffer an ischaemic stroke have a history of a preceding TIA. Of those who experience a TIA and do not receive prompt medical intervention, approximately 10\% to 15\% will go on to have a stroke within 90 days, with half of these strokes occurring within the first 48 hours. This makes the immediate post-TIA period a critical therapeutic window.

\section*{Aetiology and Pathophysiology}
The primary pathophysiology of a Transient Ischaemic Attack involves a temporary, localized reduction in cerebral, spinal, or retinal blood flow. When blood flow drops below a critical threshold, the affected nervous tissue ceases to function normally, resulting in focal neurological deficits. Because the vascular occlusion is transient---either due to spontaneous thrombolysis, fragmentation of an embolus, or collateral circulation restoration---the tissue is reperfused before irreversible cell death (infarction) occurs.

The underlying causes and risk factors for TIA are identical to those for ischaemic stroke. These can be categorized into distinct pathophysiological mechanisms:

\begin{itemize}
    \item \textbf{Large-artery atherosclerosis}: This mechanism involves the progressive accumulation of cholesterol, inflammatory cells, and fibrous tissue within the walls of major arteries supplying the brain, most commonly the carotid bifurcation and the vertebral arteries. Over time, these atherosclerotic plaques can ulcerate or rupture, exposing thrombogenic material to the bloodstream. This leads to the formation of local blood clots (thrombi) that can transiently occlude the vessel lumen, or release microemboli that travel downstream into smaller cerebral arteries, causing temporary blockages.

    \item \textbf{Cardioembolism}: Cardiogenic emboli arise from the heart chambers and are carried by the arterial system to the cerebral vasculature. The most common cause of cardioembolic TIA is atrial fibrillation (AF), a cardiac arrhythmia characterized by rapid and disorganized atrial activation. The lack of coordinated atrial contraction leads to blood stasis in the left atrial appendage, facilitating thrombus formation. Other cardiac sources of emboli include valvular heart disease (e.g., calcific aortic stenosis or mitral valve disease), prosthetic heart valves, recent myocardial infarction (which can cause mural thrombi), dilated cardiomyopathy, and patent foramen ovale (PFO)---a congenital shunt that can allow venous thrombi to bypass the lungs and enter the systemic circulation (paradoxical embolism).

    \item \textbf{Small-vessel disease (lacunar disease)}: This refers to the pathological narrowing of the small, deep penetrating arteries that branch off the major cerebral vessels to supply deep brain structures like the basal ganglia, internal capsule, and pons. This process, often called lipohyalinosis, is primarily driven by chronic hypertension and diabetes mellitus. The thickening and stiffening of the vessel walls restrict blood flow, making these areas susceptible to transient ischaemic episodes.

    \item \textbf{Hypercoagulable states and haematological disorders}: In some patients, particularly younger individuals without typical cardiovascular risk factors, a TIA may be caused by an increased propensity of the blood to clot. Inherited thrombophilias (such as factor V Leiden mutation, prothrombin gene mutation, and protein C or S deficiencies) or acquired conditions (such as antiphospholipid syndrome, polycythaemia vera, essential thrombocythaemia, and sickle cell disease) can lead to spontaneous microthrombus formation in the cerebral circulation.

    \item \textbf{Non-atherosclerotic vasculopathies}: These are less common arterial disorders that can disrupt cerebral blood flow. Examples include arterial dissection, where a tear in the inner lining of the carotid or vertebral artery allows blood to enter the vessel wall, creating a false lumen that compresses the true lumen or acts as a site for thrombus formation. Fibromuscular dysplasia (a non-inflammatory vascular disease causing abnormal cell growth in artery walls) and systemic vasculitis (inflammation of blood vessels) are other potential causes.
\end{itemize}

\section*{Clinical Features}
The clinical presentation of a Transient Ischaemic Attack is characterized by the sudden, abrupt onset of focal neurological deficits. The symptoms typically reach their maximum intensity almost instantaneously or within a few minutes. Unlike progressive neurological conditions, there is no gradual ``stuttering'' or slow evolution of symptoms. The specific deficits experienced by the patient depend entirely on the vascular territory affected and the corresponding region of the brain, spinal cord, or retina that is temporarily deprived of blood.

The clinical features of TIA are generally categorized into two main anatomical distributions: the anterior (carotid) circulation and the posterior (vertebrobasilar) circulation.

\begin{itemize}
    \item \textbf{Anterior Circulation Symptoms (Carotid Territory)}: The carotid system supplies the cerebral hemispheres, including the frontal, parietal, and temporal lobes, as well as the eyes. Anterior circulation TIAs commonly present with:
    \begin{itemize}
        \item Unilateral motor deficits, such as sudden weakness, clumsiness, or paralysis of the face, arm, and/or leg on one side of the body (hemiparesis).
        \item Unilateral sensory deficits, including numbness, tingling, or loss of sensation (hypoesthesia) affecting the face, arm, and/or leg.
        \item Speech and language disturbances, which are highly localizing to the dominant hemisphere (usually the left hemisphere). These include expressive aphasia (difficulty producing speech, word-finding difficulties), receptive aphasia (difficulty understanding spoken language), and dysarthria (slurred speech caused by weakness of the articulatory muscles).
        \item Transient monocular blindness (amaurosis fugax), which is caused by temporary ischaemia of the retina via the ophthalmic artery. Patients often describe a painless, sudden loss of vision in one eye, resembling a gray or black shade or curtain descending over their field of vision.
    \end{itemize}

    \item \textbf{Posterior Circulation Symptoms (Vertebrobasilar Territory)}: The vertebrobasilar system supplies the brainstem, cerebellum, occipital lobes, and thalamus. Posterior circulation TIAs typically present with:
    \begin{itemize}
        \item Vertigo, characterized by a false sensation of spinning or movement, which is often severe and accompanied by nausea or vomiting.
        \item Ataxia and imbalance, leading to gait instability, unsteadiness, or incoordination of the limbs.
        \item Bilateral or homonymous visual disturbances, such as complete or partial loss of vision in both eyes (e.g., homonymous hemianopia) or sudden double vision (diplopia).
        \item Dysarthria and dysphagia (difficulty swallowing) due to brainstem nuclei involvement.
        \item Alternating motor or sensory deficits, which may affect one side of the face and the opposite side of the body.
    \end{itemize}
\end{itemize}

By definition, the clinical signs of a TIA resolve completely. While historically symptoms could last up to 24 hours, the vast majority of TIA episodes are brief, with symptoms resolving within 15 to 30 minutes, and almost always within 1 hour. If symptoms persist beyond an hour, the likelihood of finding acute tissue infarction on MRI increases dramatically.

\section*{Differential Diagnosis}
Because the symptoms of a Transient Ischaemic Attack are transient and often resolved by the time the patient is evaluated by a physician, securing an accurate diagnosis can be challenging. A wide range of non-ischaemic conditions can mimic the presentation of a TIA. Distinguishing these ``TIA mimics'' from a true ischaemic event is critical, as their management and prognoses differ significantly.

The differential diagnosis of TIA includes the following conditions:

\begin{itemize}
    \item \textbf{Migraine aura}: Classical migraines are frequently accompanied by transient visual, sensory, or language disturbances (migraine aura). Unlike a TIA, which has an instantaneous onset, migraine aura typically develops gradually over 5 to 20 minutes and spreads across a sensory or visual field (sometimes referred to as the ``cortical spreading depression''). The symptoms are often ``positive'' (e.g., flashing lights, zigzag lines, tingling) rather than purely ``negative'' (e.g., vision loss, numbness). The aura is usually, but not always, followed by a throbbing headache.

    \item \textbf{Focal seizures}: A focal seizure can cause transient neurological deficits, particularly during the post-ictal phase (a phenomenon known as Todd's paresis). Todd's paresis involves a localized weakness or numbness in a limb that can last for hours after a seizure, closely resembling a TIA. However, a history of involuntary motor activity (twitching, shaking), altered consciousness, or lip-smacking during the onset points to a seizure rather than a TIA.

    \item \textbf{Hypoglycaemia}: Low blood sugar is a metabolic mimic that can produce focal neurological signs, such as hemiparesis or aphasia, especially in elderly patients or those with a history of prior stroke. It is critical to measure blood glucose levels immediately in any patient presenting with acute neurological deficits, as the symptoms resolve rapidly once glucose is administered.

    \item \textbf{Vestibular disorders}: Acute peripheral vestibular disorders, such as benign paroxysmal positional vertigo (BPPV), vestibular neuronitis, or Meniere's disease, can present with sudden, severe vertigo and balance problems. While these can mimic a posterior circulation TIA, they are distinguished by the absence of other focal brainstem or cerebellar signs (such as dysarthria, double vision, or weakness) and the presence of characteristic nystagmus.

    \item \textbf{Syncope and presyncope}: Syncope is a transient loss of consciousness caused by global cerebral hypoperfusion, often due to cardiovascular or vasovagal mechanisms. TIA, in contrast, rarely causes loss of consciousness unless there is severe, bilateral ischaemia of the brainstem reticular activating system.

    \item \textbf{Multiple Sclerosis (MS) flare}: Demyelinating plaques in MS can cause transient neurological deficits. However, these symptoms typically evolve over hours to days and persist for days or weeks, unlike the rapid onset and short duration of a TIA.

    \item \textbf{Psychogenic or conversion disorders}: Non-organic neurological disorders can present with sudden-onset weakness, sensory loss, or speech difficulty. These cases are characterized by atypical clinical findings (such as non-anatomical sensory boundaries), inconsistency on physical examination, and normal diagnostic workups.
\end{itemize}

\section*{When to Seek Medical Care}
A Transient Ischaemic Attack must never be ignored, minimized, or managed with a ``wait and see'' approach. It is a medical emergency of the highest order. Because it is impossible to determine at the onset of symptoms whether they will resolve (representing a TIA) or persist and cause permanent brain damage (representing an ischaemic stroke), every acute neurological deficit must be treated as an evolving stroke.

Public education campaigns emphasize the F.A.S.T. acronym to help individuals recognize the signs of a stroke or TIA and respond immediately:
\begin{itemize}
    \item \textbf{F - Face Drooping}: Ask the person to smile. Does one side of the face droop or is it numb?
    \item \textbf{A - Arm Weakness}: Ask the person to raise both arms. Does one arm drift downward or is it weak or numb?
    \item \textbf{S - Speech Difficulty}: Ask the person to repeat a simple sentence. Is their speech slurred, hard to understand, or are they unable to speak?
    \item \textbf{T - Time to call emergency services}: If the person shows any of these symptoms, even if they go away, call for emergency medical help immediately.
\end{itemize}

If you or someone else experiences any sudden neurological symptom---including weakness, numbness, difficulty speaking, sudden vision loss, severe dizziness, or loss of balance---you must immediately call for an ambulance.
\begin{itemize}
    \item In many countries, dial \textbf{local emergency services}.
    \item In the United States and Canada, dial \textbf{911}.
    \item In the United Kingdom, dial \textbf{999}.
    \item In Europe, dial \textbf{112}.
\end{itemize}

Do not attempt to drive yourself or the patient to the hospital. Emergency medical technicians (paramedics) can begin vital assessments, stabilize the patient, and notify the receiving hospital's stroke team. This ensures that the patient is directed to a facility with a dedicated acute stroke unit, bypassing unnecessary delays in the emergency department. Even if the symptoms resolve completely before the ambulance arrives, the patient must still be evaluated in an emergency department or specialized TIA clinic immediately.

\section*{Diagnosis and Investigations}
The primary goals of diagnostic evaluation following a suspected TIA are to confirm the clinical diagnosis, rule out mimics, identify the underlying pathophysiological mechanism, and assess the short-term risk of a subsequent stroke. Time is of the essence, and investigations should be initiated within 24 hours of symptom onset.

The diagnostic process involves a combination of clinical assessment, neuroimaging, vascular studies, and cardiac monitoring:

\begin{itemize}
    \item \textbf{Clinical Evaluation}: A thorough medical history is obtained, focusing on the onset, duration, and nature of the symptoms. A comprehensive physical and neurological examination is performed, including blood pressure measurement in both arms, auscultation of the carotid arteries for bruits, and cardiac auscultation for murmurs or irregular rhythms.

    \item \textbf{Brain Imaging}:
    \begin{itemize}
        \item \textbf{Non-contrast Head CT}: This is the initial imaging modality of choice in the emergency department. While it is rarely positive for acute ischaemia in a TIA, it is critical for ruling out intracranial haemorrhage, subdural haematomas, or intracranial masses that could mimic TIA symptoms.
        \item \textbf{Magnetic Resonance Imaging (MRI)}: MRI with Diffusion-Weighted Imaging (DWI) is the most sensitive imaging technique for detecting acute cerebral ischaemia. A positive DWI lesion indicates that tissue infarction has occurred, reclassifying the event as a minor ischaemic stroke. A negative DWI scan, in the presence of resolving focal symptoms, confirms a TIA.
    \end{itemize}

    \item \textbf{Vascular Imaging}: To identify arterial stenosis or occlusion that may have caused the TIA:
    \begin{itemize}
        \item \textbf{Carotid Duplex Ultrasound}: A non-invasive, widely available tool to assess the degree of stenosis in the extracranial carotid arteries.
        \item \textbf{CT Angiography (CTA) or MR Angiography (MRA)}: These advanced imaging modalities provide high-resolution visualization of both the extra- and intracranial vasculature, helping to detect carotid or vertebral artery stenosis, dissections, or intracranial stenosis.
    \end{itemize}

    \item \textbf{Cardiac Evaluation}: To detect potential cardioembolic sources:
    \begin{itemize}
        \item \textbf{Electrocardiogram (ECG)}: Performed immediately to look for atrial fibrillation, myocardial infarction, or left ventricular hypertrophy.
        \item \textbf{Prolonged Cardiac Monitoring}: If the initial ECG is normal but a cardioembolic source is suspected, ambulatory Holter monitoring (for 24 hours to several weeks) or insertable cardiac monitors are used to detect paroxysmal atrial fibrillation.
        \item \textbf{Echocardiography}: A transthoracic echocardiogram (TTE) or transoesophageal echocardiogram (TOE) is used to visualize the heart structure, checking for thrombi, patent foramen ovale (PFO), or valvular disease.
    \end{itemize}

    \item \textbf{Laboratory Investigations}: Routine blood tests are performed to identify metabolic risk factors. These include a complete blood count (to check for severe anaemia or thrombocytosis), coagulation studies (PT/INR, aPTT), renal function tests, blood glucose levels, HbA1c (to screen for or monitor diabetes), and a fasting lipid profile (to measure cholesterol levels).
\end{itemize}

\section*{Management}
The primary objective of TIA management is the secondary prevention of a completed ischaemic stroke. Because the risk of stroke is highest in the immediate aftermath of a TIA, management strategies must be implemented aggressively and without delay. Treatment is tailored to the specific underlying cause identified during the diagnostic workup.

Management is broadly divided into pharmacological therapy, surgical or interventional procedures, and lifestyle modifications:

\begin{itemize}
    \item \textbf{Pharmacotherapy}:
    \begin{itemize}
        \item \textbf{Antiplatelet Therapy}: For patients with non-cardioembolic TIA (e.g., those due to large-artery atherosclerosis or small-vessel disease), antiplatelets are the cornerstone of treatment.
        \begin{itemize}
            \item \textbf{Aspirin}: Initiated immediately (typically 150 mg to 300 mg loading dose, followed by 75 mg to 100 mg daily).
            \item \textbf{Dual Antiplatelet Therapy (DAPT)}: For patients with a high-risk TIA (defined by clinical scores like an ABCD2 score of $\ge$ 4 or minor stroke), the combination of aspirin and clopidogrel is initiated within 24 hours and continued for 21 to 90 days. This short-term DAPT regimen significantly reduces the risk of recurrent stroke compared to aspirin alone, without major bleeding risks.
            \item \textbf{Monotherapy}: After the initial DAPT period, patients are transitioned to long-term single antiplatelet therapy, typically with clopidogrel (75 mg daily) or low-dose aspirin.
        \end{itemize}
        \item \textbf{Anticoagulation}: For patients with cardioembolic TIA, particularly those with atrial fibrillation, long-term anticoagulation is required. Direct Oral Anticoagulants (DOACs) such as apixaban, rivaroxaban, dabigatran, or edoxaban are preferred over warfarin because they offer similar efficacy with a lower risk of life-threatening bleeding (specifically intracranial haemorrhage) and do not require routine blood monitoring.
        \item \textbf{Lipid-Lowering Therapy}: High-intensity statin therapy (such as atorvastatin 80 mg daily) is prescribed for all patients with a TIA of atherosclerotic origin, regardless of their baseline cholesterol levels. Statins lower LDL cholesterol, stabilize atherosclerotic plaques, and have anti-inflammatory effects. The target LDL level is typically less than 1.8 mmol/L (70 mg/dL) or a 50\% reduction from baseline.
        \item \textbf{Antihypertensive Therapy}: Long-term blood pressure control is vital. Antihypertensive therapy (commonly using ACE inhibitors, angiotensin receptor blockers, or thiazide diuretics) is initiated or optimized to achieve a target blood pressure of less than 130/80 mmHg.
    \end{itemize}

    \item \textbf{Surgical and Interventional Procedures}:
    \begin{itemize}
        \item \textbf{Carotid Endarterectomy (CEA)}: For patients with moderate-to-severe (50\% to 99\%) stenosis of the symptomatic internal carotid artery, surgical removal of the plaque (CEA) is highly effective. To achieve maximum benefit, the procedure should be performed within 14 days of the TIA.
        \item \textbf{Carotid Artery Stenting (CAS)}: A less invasive alternative to CEA, where a metal stent is placed inside the carotid artery to keep it open. CAS is typically reserved for patients who are at high risk for surgical complications, have anatomically inaccessible lesions, or have radiation-induced stenosis.
    \end{itemize}

    \item \textbf{Therapies and Support}: While TIA symptoms resolve, some patients may experience transient speech or motor difficulties. A brief consultation with physical, occupational, or speech therapy can help ensure complete recovery and provide guidance on safe return to daily activities.
\end{itemize}

\section*{Complications}
While a Transient Ischaemic Attack itself resolves without causing permanent direct damage, the primary complication is the development of a subsequent, disabling, or fatal ischaemic stroke. A TIA is a warning sign that the cerebrovascular system is compromised.

The potential complications and risks associated with TIA and its treatments include:

\begin{itemize}
    \item \textbf{Subsequent Ischaemic Stroke}: This is the most feared outcome. If left untreated, up to 15\% of TIA patients will suffer a stroke within 90 days, with the highest risk concentrating in the first 48 hours. A completed stroke can result in permanent neurological deficits, such as hemiplegia, aphasia, cognitive impairment, or death.

    \item \textbf{Cognitive Decline and Vascular Dementia}: Repeated episodes of subclinical cerebral ischaemia, even if they do not cause noticeable symptoms, can accumulate damage in the brain's white matter. Over time, this microvascular damage can lead to vascular cognitive impairment or vascular dementia, affecting memory, planning, and executive function.

    \item \textbf{Bleeding Complications from Pharmacotherapy}: The medications used to prevent stroke---antiplatelets and anticoagulants---compromise the body's clotting ability. Consequently, patients face an increased risk of bleeding. This ranges from minor bruising and epistaxis (nosebleeds) to major, life-threatening events such as gastrointestinal bleeding or haemorrhagic stroke (bleeding into the brain).

    \item \textbf{Procedural Risks}: Surgical interventions like carotid endarterectomy (CEA) or carotid artery stenting (CAS) carry inherent risks. These include perioperative stroke (caused by plaque debris breaking off during the procedure), myocardial infarction, local wound infection, haematoma, damage to cranial nerves (affecting swallowing or voice), and cerebral hyperperfusion syndrome (a rare but serious condition causing headache, seizures, or haemorrhage due to sudden restoration of normal blood flow to a chronically hypoperfused brain).

    \item \textbf{Psychological distress}: Many patients experience significant anxiety, depression, and post-traumatic stress symptoms following a TIA. The sudden nature of the event, coupled with the knowledge that it is a ``warning stroke,'' can lead to a persistent fear of recurrence, impairing daily functioning and quality of life.
\end{itemize}

\section*{Prognosis and Outcomes}
The prognosis of an individual who has experienced a Transient Ischaemic Attack is highly variable and depends on several factors, including the underlying cause, the presence of comorbid medical conditions, the patient's age, and---most importantly---the speed and efficacy of medical intervention.

In the absence of prompt evaluation and treatment, the prognosis is concerning due to the high risk of a subsequent stroke. Historically, approximately 10\% to 15\% of patients experienced a stroke within three months of a TIA, with half of those strokes occurring within 48 hours. To help clinicians stratify this risk and identify patients who require immediate hospitalization, clinical prediction tools like the \textbf{ABCD2 score} were developed. The ABCD2 score assigns points based on:
\begin{itemize}
    \item \textbf{A - Age}: 1 point for age $\ge$ 60 years.
    \item \textbf{B - Blood Pressure}: 1 point for systolic BP $\ge$ 140 mmHg or diastolic BP $\ge$ 90 mmHg.
    \item \textbf{C - Clinical Features}: 2 points for unilateral weakness; 1 point for speech disturbance without weakness.
    \item \textbf{D - Duration}: 2 points for symptoms lasting $\ge$ 60 minutes; 1 point for symptoms lasting 10 to 59 minutes.
    \item \textbf{D - Diabetes}: 1 point for pre-existing history of diabetes.
\end{itemize}
A higher ABCD2 score (5 to 7) indicates a high risk of stroke in the subsequent 48 hours, warranting urgent inpatient admission and rapid workup.

However, with modern, aggressive management strategies, the prognosis has improved dramatically. Studies have demonstrated that rapid evaluation in a specialized TIA clinic or emergency department, followed by immediate initiation of dual antiplatelet therapy, high-intensity statins, blood pressure optimization, and timely carotid surgery (when indicated), reduces the risk of subsequent stroke by up to 80\%. For patients who receive timely and appropriate care, the long-term prognosis is excellent, and they can expect to return to their baseline level of functioning.

\section*{Prevention and Self-Care}
Preventing a recurrence of TIA and avoiding a future completed stroke requires a comprehensive approach that combines adherence to prescribed medical therapies with active self-care and lifestyle modification. Because many cardiovascular risk factors are modifiable, patients play a central role in managing their own prognosis.

Key self-care and prevention measures include:

\begin{itemize}
    \item \textbf{Blood Pressure Control}: Hypertension is the single most important modifiable risk factor for stroke. Patients should monitor their blood pressure regularly at home using a validated upper-arm monitor, keep a log of the readings, and aim for a target of less than 130/80 mmHg. Strict compliance with prescribed antihypertensive medications is essential.

    \item \textbf{Healthy Diet}: Adopting a Mediterranean-style diet has been shown to reduce cardiovascular risk. This diet emphasizes high consumption of vegetables, fruits, whole grains, legumes, nuts, and olive oil; moderate consumption of fish and poultry; and low consumption of red meat, processed foods, and sweets. Reducing dietary sodium (salt) to less than 2,000 mg per day helps lower blood pressure.

    \item \textbf{Regular Physical Activity}: Engaging in regular exercise improves cardiovascular health, helps control weight, and lowers blood pressure. Unless medically contraindicated, patients should aim for at least 150 minutes of moderate-intensity aerobic physical activity (such as brisk walking, cycling, or swimming) per week, distributed over several days.

    \item \textbf{Smoking Cessation}: Tobacco use significantly accelerates atherosclerosis, damages blood vessel linings, and increases the blood's propensity to clot. Quitting smoking is one of the most impactful changes a patient can make. Healthcare providers can offer support through nicotine replacement therapies, counselling, and prescription medications.

    \item \textbf{Limiting Alcohol Consumption}: High alcohol intake is associated with hypertension, atrial fibrillation, and an increased risk of stroke. Men should limit their consumption to no more than two standard drinks per day, and women to no more than one standard drink per day.

    \item \textbf{Weight Management}: Achieving and maintaining a healthy body weight (defined as a body mass index between 18.5 and 24.9 kg/m$^2$) helps control blood pressure, cholesterol, and diabetes risk.

    \item \textbf{Medication Adherence}: Patients must take all prescribed medications---including antiplatelets, anticoagulants, statins, and blood pressure medications---exactly as directed. They should never stop taking these medications or alter the dose without consulting their physician, even if they feel well, as the protective benefits are only maintained with continuous use.

    \item \textbf{Managing Diabetes}: For patients with diabetes, maintaining strict glycaemic control (aiming for an HbA1c level recommended by their physician, typically below 7.0\%) through diet, exercise, and medication is crucial for protecting the cerebral microvasculature.
\end{itemize}

\section*{Sources and Further Reading}
For further information, support, and clinical guidelines on Transient Ischaemic Attack, the following authoritative resources are recommended:

\begin{itemize}
    \item \textbf{Stroke Foundation (Australia)}: An independent, not-for-profit organisation dedicated to stroke prevention, support for survivors, and clinical education. Their website offers comprehensive guides on recognizing TIA and stroke symptoms, lifestyle management, and recovery. Available at \texttt{https://strokefoundation.org.au}.

    \item \textbf{American Stroke Association (ASA)}: A division of the American Heart Association, the ASA provides patient education materials, scientific statements, and clinical guidelines for stroke and TIA prevention and treatment. Available at \texttt{https://www.stroke.org}.

    \item \textbf{World Stroke Organization (WSO)}: The leading global body in the fight against stroke, representing stroke support organisations, national societies, and individual clinicians. The WSO publishes global guidelines and educational campaigns. Available at \texttt{https://www.world-stroke.org}.

    \item \textbf{National Institute for Health and Care Excellence (NICE)}: The UK-based independent organisation providing evidence-based national guidance and advice on health and social care. NICE guidelines on stroke and TIA cover initial assessment, referral times, and secondary prevention. Available at \texttt{https://www.nice.org.uk}.

    \item \textbf{Mayo Clinic}: A world-renowned academic medical center providing detailed, patient-focused articles on the symptoms, causes, diagnosis, and management of TIA. Available at \texttt{https://www.mayoclinic.org}.

    \item \textbf{National Institute of Neurological Disorders and Stroke (NINDS)}: Part of the United States National Institutes of Health (NIH), NINDS conducts and supports research on brain and nervous system disorders. Their resources provide scientific overviews of TIA and current clinical trials. Available at \texttt{https://www.ninds.nih.gov}.
\end{itemize}